\documentclass[11pt]{article}
\usepackage{amsmath,amssymb}
\usepackage[margin=1in]{geometry}
\title{Weekly research summary}
\author{Daniel}
\date{7 October 2026}
\begin{document}
\maketitle
\section*{Grünbaum project}
This week I concentrated on understanding the deformation-theoretic input to the current \(v2\) proof rather than treating the \(T^1\) and \(T^2\) computations as black boxes.

\begin{itemize}
\item I reconstructed the cotangent complex from Kähler differentials. For \(R=P/I\), the conormal sequence is
$$
I/I^2\longrightarrow\Omega^1_{P/k}\otimes_PR
\longrightarrow\Omega^1_{R/k}\longrightarrow0.
$$
The general construction resolves \(R\) by polynomial algebras \(P_\bullet\to R\) and sets
$$
L_{R/k}=\Omega^1_{P_\bullet/k}\otimes_{P_\bullet}R.
$$

\item I rebuilt first-order deformations from the dual numbers \(D=k[\varepsilon]/(\varepsilon^2)\). A first-order deformation is a flat \(D\)-algebra \(\widetilde R\) with
$$
\widetilde R/\varepsilon\widetilde R\simeq R.
$$
Flatness gives \(\varepsilon\widetilde R\simeq R\), and the kernel is square-zero.

\item After choosing a \(k\)-linear section, the failure to preserve multiplication is measured by \(\mu:R\times R\to R\). Associativity gives
$$
a\mu(b,c)-\mu(ab,c)+\mu(a,bc)-\mu(a,b)c=0,
$$
while changing the section by \(h:R\to R\) changes \(\mu\) by
$$
a h(b)-h(ab)+h(a)b.
$$
Thus the concrete deformation data is compatible \(\mu\)'s modulo changes from \(h\). This is the Hochschild cocycle/coboundary picture for associative deformations.

\item I reviewed why equivalence classes of module extensions
$$
0\longrightarrow A\longrightarrow E\longrightarrow B\longrightarrow0
$$
are \(\operatorname{Ext}^1_R(B,A)\), including the projective-resolution lifting argument and the splitting lemma.

\item I defined cotangent cohomology
$$
T^i_{R/k}:=\operatorname{Ext}^i_R(L_{R/k},R).
$$
I am now proving the bridge from the concrete \(\mu/h\) calculation to the theorem that \(T^1\) classifies first-order commutative algebra deformations. The current step is the algebra analogue of the ordinary Ext argument: for a square-zero extension
$$
0\longrightarrow M\longrightarrow\widetilde R\longrightarrow R\longrightarrow0,
$$
the map \(P_0\to R\) from the degree-zero polynomial algebra in \(P_\bullet\to R\) lifts to \(P_0\to\widetilde R\) by lifting its polynomial generators.
\end{itemize}

\section*{Where I am in the paper}
This work is immediately upstream of the explicit deformation-theory calculations in the Grünbaum \(v2\) proof. The computational input includes
$$
\dim T^1_0=53,\qquad \dim T^2_0=27.
$$
I am currently making precise why \(T^1_0\) is the graded first-order deformation space and why \(T^2_0\) is the relevant obstruction space. After this, I return to the paper-specific multigraded description of \(T^2_0\) (including the Altmann--Christophersen pieces such as \(a^2/(cg)\)), and then to the symmetry and Maurer--Cartan/Bianchi all-orders lifting argument.

\section*{Next step}
Finish the proof that square-zero algebra extensions by an \(R\)-module \(M\) are classified by
$$
\operatorname{Ext}^1_R(L_{R/k},M),
$$
specialize to \(M=R\), and reconnect the intrinsic picture to the explicit \(T^1_0,T^2_0\) data in the Grünbaum manuscript.
\end{document}
